\chapter{Funkcije}

\section{Pojam funkcionalne zavisnosti}

Podaci o jednoj (izmi\v sljenoj!) grupi ljudi su predstavljeni slede\'com tabelom:
\begin{center}
  \begin{tabular}{@{}rlllr@{}}
  \hline
  JMBG & Prezime & Ime & Adresa & Datum ro\dj. \\
  \hline
  3112965870015 & Petrovi\'c & Petar & Bul.\ mimoze 21, Novigrad & 31.12.1965.\\
  3011967875021 & Zori\'c & Zora & Bul.\ mimoze 21, Novigrad & 30.11.1967.\\
  1205995845039 & Mari\'c & Mara & Avenija ru\v ze, Maligrad & 12.5.1995.\\
  2808985830029 & Markovi\'c & Petar & Ul.\ razli\v cka 25, Veligrad &  28.8.1985.\\
  2808985810034 & Petrovi\'c & Rade & Bulkin put bb, Maligrad &  28.8.1985.
  \end{tabular}
\end{center}
Jasno je da je svaka od ovih osoba jednozna\v cno odre\dj ena paketom podataka koji sadr\v zi
prezime, ime, adresu i datum ro\dj enja. Ali nijedan od tih podataka za sebe nije dovoljan da
odredi osobu (postoje dve osobe sa istim prezimenom, dve osobe sa istim imenom, dve osobe sa istom adresom
i dve osobe sa istim datumom ro\dj enja). Da bi se obrada podataka pojednostavila i ubrzala
veoma je va\v zno da u svakoj grupi podataka postoji atribut koji svojim vrednostima
jednozna\v cno odre\dj uje podatke. U Republici Srbiji je
to broj koji se zove Jedinstveni mati\v cni broj gra\dj anina (JMBG).

Funkcionalna zavisnost je veza me\dj u atributima nekog skupa podataka koja se ogleda u tome da
vrednost jednog atributa jednozna\v cno odre\dj uje ceo podatak, kao \v sto JMBG jednozna\v cno odre\dj uje
ostale podatke o osobi. Dakle, funkcionalna zavisnost je relacija sa posebnim osobinama.

Formalno, ako su $A$ i $B$ dva skupa i ako je $\rho \subseteq A \times B$ neka relacija, ka\v zemo da
ona opisuje \emph{funkcionalnu zavisnost} izme\dj u ta dva skupa ako
za svaki $a \in A$ postoji ta\v cno jedan $b \in B$ tako da je $(a, b) \in \rho$. Dakle, svaki element skupa
$A$ odre\dj uje ta\v cno jedan element skupa $B$.

U matematici se \v ce\v s\'ce umesto ``funkcionalna zavisnost'' ka\v ze ``funkcija'' i umesto malih gr\v ckih
slova (koja smo rezervisali za relacije) kao oznaku koristimo mala latini\v cna slova. Tako, umesto
``$\rho \subseteq A \times B$ je funkcionalna zavisnost'' radije \'cemo re\'ci
``$f : A \to B$ je funkcija'', a umesto ``$(a, b) \in \rho$'' pi\v semo ``$f(a) = b$''.
Uz sve navedne konvencije definicija funkcije sada glasi ovako:

\begin{quote}
  Funkcija iz skupa $A$ u skup $B$ je svaka relacija $f \subseteq A \times B$ koja ima slede\'cu osobinu:
  za svako $a \in A$ postoji ta\v cno jedno $b \in B$ takvo da je $f(a) = b$. Pi\v semo $f : A \to B$
  ili $A \overset{f}{\rightarrow} B$.
  Skup $A$ je \emph{domen} funkcije $f$, a skup $B$ njen \emph{kodomen}.
\end{quote}

\begin{EX}
  Funkciju mo\v zemo da opi\v semo na razne na\v cine. Neka je
  $$
    \Phi = \big\{\bigtriangleup, \Diamond, \square, \bigcirc, \pentagon, \hexagon, \triangleleft, \bowtie\big\}
  $$
  skup geometrijskih figura, a $t : \Phi \to \ZZ$ funkcija koja figuri dodeljuje broj temena.
  Funkciju $t$ mo\v zemo da opi\v semo prosto navo\dj enjem svih mogu\'cnosti:
  $$
    \begin{array}{llll}
      t\big(\bigtriangleup\big) = 3, & t(\Diamond) = 4, & t(\square) = 4, & t(\bigcirc) = 0,\\
      t(\pentagon) = 5, & t(\hexagon) = 6, & t(\triangleleft) = 3, & t(\bowtie) = 5,
    \end{array}
  $$
  ili tabelom:
  $$
    \begin{array}{c|cccccccc}
      \phi    & \bigtriangleup & \Diamond & \square & \bigcirc & \pentagon & \hexagon & \triangleleft & \bowtie\\
      \hline
      t(\phi) &  3             & 4        & 4       & 0        & 5         & 6        & 3             & 5
    \end{array}
  $$
  ili kao skup ure\dj enih parova (relacija):
  $$
    t = \big\{\big(\bigtriangleup, 3\big),\; (\Diamond, 4),\; (\square, 4),\; (\bigcirc, 0),\; (\pentagon, 5),\;
                  (\hexagon, 6),\; (\triangleleft, 3),\; (\bowtie, 5) \big\}
  $$
  ili Venovim dijagramom:
  \begin{center}
    \input ft08-f1.pgf
  \end{center}
\end{EX}

\begin{EX}
  Ako su domen i kodomen funkcije veoma veliki, nije uobi\v cajeno da se funkcija predstavlja na neki od na\v cina koji
  su opisani u prethodnom primeru, ve\'c se koriste druge konvencije. Na primer, jedan na\v cin da se opi\v se
  funkcija je da se zada formula koja odre\dj uje na\v cin na koji se elementi njenog domena preslikavaju na elemente kodomena.
  Na primer, neka je $f : \NN \to \NN$ funkcija zadata ovako:
  $$
    f(n) = \frac{n(n + 1)(2n + 1)}{6}
  $$
  (lako se pokazuje da je broj $n(n + 1)(2n + 1)$ uvek deljiv sa 6, tako da ova funkcija zaista preslikava prirodne brojeve
  na prirodne brojeve). Kao drugi primer navodimo funkciju $| \cdot | : \RR \to \RR$ koja je zadata ovako
  (apsolutna vrednost realnog broja):
  $$
    |x| = \begin{cases}
      -x, & x < 0,\\
      x,  & x \ge 0.
    \end{cases}
  $$
  Kao poslednji primer navodimo funkciju $\mathrm{sgn} : \RR \to \RR$ koja odre\dj uje znak broja (funkcija signum):
  $$
    \mathrm{sgn} \; x = \begin{cases}
      -1, & x < 0,\\
       0, & x = 0,\\
       1, & x > 0.
    \end{cases}
  $$
\end{EX}

\begin{EX}
  Realni koren \emph{nije funkcija $\RR \to \RR$} jer ne mo\v ze da se primeni na sve elemente domena (realni koren
  nije definisan za negativne realne brojeve). Ako sa $\RR_{\ge 0} = \{x \in \RR : x \ge 0\}$ ozna\v cimo skup svih
  nenegativnih relasnih brojeva, onda je $\sqrt{\mathstrut\quad} : \RR_{\ge 0} \to \RR$ funkcija.
\end{EX}

\begin{EX}
  Ponekad je funkciju pogodno zadati grafikom. Recimo, grafik funkcije
  $\sin : \RR \to \RR$ izgleda ovako:
  \begin{center}
    \input ft08-f2.pgf
  \end{center}
\end{EX}

\begin{EX}\label{ft08.ex.strlen}
  Neka je $A$ alfabet (kona\v can neprazan skup \v cije elemente zovemo slova), i neka je $A^*$ skup svih
  re\v ci nad alfabetom $A$. Tada je sa $\mathrm{len} : A^* \to \NN \union \{0\}$ ozna\v cena funkcija koja
  svakoj re\v ci pridru\v zuje njenu du\v zinu (broj slova u toj re\v ci). Na primer, za $A = \{a, b, c\}$
  je $\mathrm{len}(aabab) = 5$ i $\mathrm{len}(\epsilon) = 0$.
  
  U programskim jezicima se ova funkcija ozna\v cava sa \verb`length`, \verb`len` ili tako nekako. Na primer,
  u programskom jeziku Java tip \verb`String` opisuje kona\v cne nizove simbola, a du\v zina stringa se
  dobija pozivom funkcije \verb`length()` ovako:
{\small
\begin{verbatim}
    Scanner ulaz = new Scanner(System.in);
    String s = ulaz.nextLine();
    System.out.println(s.length());
\end{verbatim}
}
\end{EX}

\begin{EX}\label{ft08.ex.stanje-mem}
  Stanje promenljivih u memoriji ra\v cunara mo\v zemo da opi\v semo kao funkciju koja imenu promenljive
  pridru\v zuje njenu vrednost. Neka je $A$ alfabet i neka je
  $$
    A^+ = A \cup A^2 \cup A^3 \cup \ldots = A^* \setminus \{\epsilon\}
  $$
  skup svih \emph{nepraznih re\v ci} nad alfabetom $A$. Posmatrajmo pojednostavljen model u kome promenljive mogu da imaju
  samo celobrojne vrednosti. Tada je \emph{stanje memorije} svaka funkcija $S : A^+ \to \ZZ$. Za $W \in A^+$ sa $S(W)$ je
  ozna\v cena vrednost promenljive~$W$. Na primer, neka je $A = \{a, \ldots, z\}$ skup svih malih slova engleskog alfabeta.
  Jedno stanje memorije je opisano funkcijom $\mathrm{INIT}_0$ koja sve promenljive inicijalizuje na vrednost~0:
  $$
    \begin{array}{c|cccc}
      W & a & b & c & \ldots\\
      \hline
      \mathrm{INIT}_0(W) & 0 & 0 & 0 & \ldots
    \end{array}
  $$
  Neko drugo stanje memorije $T : A^+ \to \ZZ$ mo\v ze da izgleda, na primer, ovako:
  $$
    \begin{array}{c|cccc}
      W & a & b & c & \ldots\\
      \hline
      T(W) & -5 & 0 & 7 & \ldots
    \end{array}
  $$
  Neka je $\mathrm{MEM}$ skup svih mogu\'cih stanja memorije. Svaka naredba programskog jezika menja stanje memorije,
  i zato svakoj naredbi programskog jezika odgovara neka funkcija $\mathrm{MEM} \to \mathrm{MEM}$. Na primer, naredbi
  $a = 1$ odgovara funkcija\break
  $N_{a = 1} : \mathrm{MEM} \to \mathrm{MEM}$ koja stanje $S$ preslikava u stanje $S'$ koje se od
  stanja $S$ razlikuje samo to tome \v sto je $S'(a) = 1$. Na primer za stanje $T$ koje je prikazano gore
  je $N_{a = 1}(T) = T'$, gde je stanje $T'$ dato sa:
  $$
    \begin{array}{c|cccc}
      W & a & b & c & \ldots\\
      \hline
      T'(W) & 1 & 0 & 7 & \ldots
    \end{array}
  $$
  Sli\v cno tome, naredbi $c = c + 3$ odgovara funkcija $N_{c = c + 3} : \mathrm{MEM} \to \mathrm{MEM}$
  koja stanje $S$ preslikava u stanje $S'$ koje se od stanja $S$ razlikuje samo to tome \v sto je $S'(c) = S(c) + 3$.
  Na primer za stanje $T'$ koje je prikazano gore je $N_{c = c + 3}(T') = T''$, gde je stanje $T''$ dato sa:
  $$
    \begin{array}{c|cccc}
      W & a & b & c & \ldots\\
      \hline
      T''(W) & 1 & 0 & 10 & \ldots
    \end{array}
  $$
  Ovom primeru \'cemo se vratiti kada budemo diskutovali kompoziciju finkcija.
\end{EX}

\clearpage

\noindent
Funkcije $f : A \to B$ i $g : C \to D$ su \emph{jednake} ako
\begin{itemize}\itemsep -1.5pt
\item
  $A = C$ i $B = D$ (imaju isti domen i kodomen), i
\item
  $f(x) = g(x)$ za sve $x \in A = C$ (pona\v saju se na isti na\v cin).
\end{itemize}

\begin{EX}
  Neka je $f : \NN \to \RR$ funkcija data sa $f(x) = x^2$ i neka je $g : \RR \to \RR$ funkcija data sa $g(x) = x^2$.
  Iako ove dve funkcije ``rade isti posao'' -- kvadriraju broj -- one nisu jednake jer imaju razli\v cite domene.
\end{EX}

\begin{EX}
  Neka je $f : \RR \to \RR$ funkcija data sa $f(x) = \sin(x)$, a $g : \RR \to [0,1]$ funkcija data sa $g(x) = \sin x$.
  I ove dve funkcije ``rade isti posao'' -- obe ra\v cunaju sinus broja -- ali one nisu jednake jer imaju razli\v cite kodomene.
\end{EX}

\begin{EX}
  Neka je $f : \RR \to \RR$ funkcija data sa $f(x) = 1$, a $g : \RR \to \RR$ funkcija data sa $g(x) = \sin^2 x + \cos^2 x$.
  Ovde dve funkcije su jednake jer imaju iste domene i kodomene i zato \v sto je $f(x) = g(x) = 1$ za sve $x \in \RR$.
\end{EX}

\section{Kompozicija funkcija}

\noindent
\begin{minipage}{3in}
  \hspace*{1.5em}Ako su $f : A \to B$ i $g : B \to C$ funkcije gde je kodomen prve jednak domenu druge, one se mogu komponovati
  kako bi se dobila nova funkcija\break
  $g \circ f : A \to C$ koja je definisana tako da je
  $$
    (g \circ f)(x) = g(f(x)), \text{\ \ za sve\ \ } x \in A.
  $$
  \v Citamo: ``$g$ kru\v zi\'c $f$''.
\end{minipage}
\hfill
\begin{minipage}{2.5in}
  \begin{center}
    \input ft08-f3.pgf
  \end{center}
\end{minipage}

\medskip

Funkcija $g \circ f$ deluje tako da se na $x$ prvo primeni $f$, pa se na rezultat $f(x)$ primenjuje funkcija $g$. Zato se
$g \circ f$ ponekad ozna\v cava i sa $f \mathbin; g$ i \v cita ``prvo $f$, onda~$g$''. Dakle,
$$
 (f \mathbin; g)(x) = (g \circ f)(x) = g(f(x)).
$$

\begin{EX}
  Neka su $f : \RR \to \RR_{\ge 0}$ i $g : \RR_{\ge 0} \to \RR$ realne funkcije date sa
  $f(x) = x^2$ i $g(x) = \sqrt x$, gde smo sa $\RR_{\ge 0}$ ozna\v cili skup svih nenegativnih realnih brojeva.
  Tada je $g \circ f : \RR \to \RR$ funkcija data sa
  $$
    (g \circ f)(x) = g(f(x)) = g(x^2) = \sqrt{x^2} = |x|,
  $$
  dok je $f \circ g : \RR_{\ge 0} \to \RR_{\ge 0}$ funkcija data sa
  $$
    (f \circ g)(x) = f(g(x)) = f(\sqrt x) = (\sqrt{x})^2 = x.
  $$
\end{EX}

\begin{EX}
  Neka je $A = \{0, 1\}$ alfabet, neka je $\mathrm{len} : A^* \to \NN \cup \{0\}$ funkcija koja re\v ci dodeljuje njenu
  du\v zinu (v.~Primer~\ref{ft08.ex.strlen}), a $\beta : \NN \to A^*$ funkcija koja svakom prirodnom broju dodeljuje
  njegov zapis u binarnom sistemu (na primer, $\beta(1) = 1$, $\beta(9) = 1001$ itd). Tada je
  $\mathrm{len} \circ \beta : \NN \to \NN \union \{0\}$ funkcija koja ra\v cuna du\v zinu binarnog zapisa prirodnog broja.
  Na primer,
  $$
    (\mathrm{len} \circ \beta)(9) = \mathrm{len}(\beta(9)) = \mathrm{len}(1001) = 4.
  $$
\end{EX}

\begin{EX}
  Vratimo se na Primer~\ref{ft08.ex.stanje-mem} u kome smo opisali jedan jednostavan model stanja memorije.
  Da se podsetimo, sa $A$ smo ozna\v cili alfabet. Neprazne re\v ci nad tim alfabetom odgovaraju imenima
  promenljivih za koje smo pretpostavili da mogu da imaju samo celobrojne vrednosti.
  Stanje promenljivih u memoriji ra\v cunara mo\v zemo da opi\v semo kao funkciju $S : A^+ \to \ZZ$
  koja imenu promenljive pridru\v zuje njenu vrednost. Sa $\mathrm{MEM}$ smo ozna\v cili skup svih mogu\'cih stanja memorije.
  
  Svaka naredba programskog jezika menja stanje memorije,
  i zato svakoj naredbi programskog jezika odgovara neka funkcija $\mathrm{MEM} \to \mathrm{MEM}$.
  Izvr\v savanje dve naredbe jedne za drugom ta\v cno odgovara kompoziciji odgovaraju\'cih funkcija.
  Na primer, neka je $T$ slede\'ce stanje memorije:
  $$
    \begin{array}{c|cccc}
      W & a & b & c & \ldots\\
      \hline
      T(W) & 1 & 0 & 7 & \ldots
    \end{array}
  $$
  Niz naredbi $a = 9;\; a = a + 3$ odgovara kompoziciji $N_{a = 9} \mathbin; N_{a = a + 3} : \mathrm{MEM} \to \mathrm{MEM}$
  koja na stanje $T$ deluje ovako
  $$
    (N_{a = 9} \mathbin; N_{a = a + 3})(T) = N_{a = a + 3}(N_{a = 9}(T))
  $$
  (prvo $N_{a = 9}$, pa posle toga $N_{a = a + 3}$, zato \v sto su naredbe navedene tim redom). Navedeni niz funkcija menja
  stanje memorije kako i o\v cekujemo:
  \begin{center}
    \input ft08-f4.pgf
  \end{center}
\end{EX}

\clearpage

\begin{THM}
  Neka su $f : A \to B$, $g : B \to C$ i $h : C \to D$ tri funkcije. Tada je $(f \circ g) \circ h = f \circ (g \circ h)$.
\end{THM}
\begin{proof}[Dokaz]
  Jasno je da i funkcija $(f \circ g) \circ h$ i funkcija $f \circ (g \circ h)$ imaju domen $A$ i kodomen $D$, tako da jo\v s
  samo treba dokazati da je $((f \circ g) \circ h)(x) = (f \circ (g \circ h))(x)$, za svako $x \in A$.
  Mada na prvi pogled deluje komplikovano, ovo se dokazuje pravolinijski. S jedne strane je
  $$
    ((f \circ g) \circ h)(x) = (f \circ g)(h(x)) = f(g(h(x))).
  $$
  S druge strane,
  $$
    (f \circ (g \circ h))(x) = f((g \circ h)(x)) = f(g(h(x))).
  $$
  Dakle, za svako $x \in A$ je $((f \circ g) \circ h)(x) = f(g(h(x))) = (f \circ (g \circ h))(x)$,
  \v cime je dokaz zavr\v sen.
\end{proof}



\section{Injektivne, sirjektivne i bijektivne funkcije}

Funkcija $f : A \to B$ je \emph{injektivna} (ili ``1-1'') ako za sve $a_1, a_2 \in A$ imamo
$$
  a_1 \ne a_1 \Rightarrow f(a_1) \ne f(a_2).
$$
\v Cesto se injektivna funkcija defini\v se ovako: funkcija
$f : A \to B$ je \emph{injektivna} (ili ``1-1'') ako za sve $a_1, a_2 \in A$ imamo
$$
  f(a_1) = f(a_2) \Rightarrow a_1 = a_2.
$$
Ove dve definicije su ekvivalentne zato je drugi zahtev kontrapozicija prvog, a kontrapozitivna
tvr\dj enja su ekvivalentna. Dakle, svejedno je koju \'cemo definiciju usvojiti.

\begin{EX}
  Neka su elementi skupa $P$ nazivi proizvoda u jednoj prodavnici. \v Sifra proizvoda je prirodan broj
  dodeljen nazivu proizvoda. Dakle, \v sifra proizvoda je funkcija
  $$
    \hbox{\v sifra} : P \to \NN.
  $$
  Jasno je da ne mora svaki prirodan broj biti upotrebljen kao \v sifra nekog proizvoda
  (na kraju krajeva, skup prirodnih brojeva je beskona\v can, a skup proizvoda kona\v can).
  Da bi \v sifrovanje bilo smisleno moramo razli\v citim proizvodima dodeliti razli\v cite \v sifre:
  za svaka dva proizvoda $p_1, p_2 \in P$ treba da va\v zi
  $$
    p_1 \ne p_2 \Rightarrow \hbox{\v sifra}(p_1) \ne \hbox{\v sifra}(p_2).
  $$
  \begin{center}
    \input ft08-f5.pgf
  \end{center}
  Dakle, svako smisleno \v sifrovanje je injektivna ili ``1-1'' funkcija (oznaka ``1-1''
  poti\v ce od ove slike: jedan proizvod --- jedna \v sifra).
\end{EX}

\begin{EX}
  U programskom jeziku Java imamo tip \verb`int` koji opisuje deo skupa $\ZZ$, i tip \verb`double` koji
  opisuje deo skupa $\RR$. Iako je $\ZZ \subseteq \RR$, u programskom jeziku Java
  $$
    \hbox{\texttt{int}} \not\subseteq \hbox{\texttt{double}}
  $$
  zato \v sto se ma\v sinska reprezentacija celobrojnih vrednosti su\v stinski razlikuje od ma\v sinske reprezentacije
  decimalnih vrednosti. \v Stavi\v se,
  $$
    \hbox{\texttt{int}} \cap \hbox{\texttt{double}} = \0.
  $$
  Me\dj utim, u programskom jeziku Java postoji funkcija
  $$
    \hbox{\texttt{(double)}} : \hbox{\texttt{int}} \to \hbox{\texttt{double}}
  $$
  koja celobrojnu vrednost transformi\v se u odgovaraju\'cu decimalnu vrednost. Jasno je da je funkcija
  \verb`(double)` injektivna jer razli\v cite cele brojeve preslikava na razli\v cite decimalne brojeve.
  Ova konstrukcija se u programskom jeziku Java zove \emph{type cast} i koristi se, na primer, ovako:
  {\small
\begin{verbatim}
    int k = 9;
    double x = (double)k;
\end{verbatim}
  }
\end{EX}

\bigskip

Funkcija $f : A \to B$ je \emph{sirjektivna} (ili ``na'') ako za svako $b \in B$ postoji
$a \in A$ takvo da je $f(a) = b$.

\begin{EX}
  Jedna \v skola ide na ekskurziju i za to je anga\v zovala nekoliko autobusa.
  Neka je $U$ skup koji opisuje u\v cenike \v skole (recimo da je svaki u\v cenik predstavljen svojim JMBG brojem)
  i neka je $B$ skup koji opisuje autobuse (recimo da je svaki autobus predstavljen svojim registarskim brojem).
  Raspodela u\v cenika po autobusima se tada mo\v ze predstaviti kao funkcija
  $$
    \mathit{bus} : U \to B
  $$
  koja svakom u\v ceniku $u \in U$ dodeljuje autobus $\mathit{bus}(u)$. Da bi ekskurzija bila rentabilna
  jedan od osnovnih zahteva je da u ovoj raspodeli nema praznih autobusa\footnote{realnost je, naravno, jo\v s zahtevnija}.
  To zna\v ci da
  \begin{quote}
    za svaki $b \in B$ postoji bar jedan $u \in U$ takav da je $\mathit{bus}(u) = b$.
  \end{quote}
  \begin{center}
    \input ft08-f6.pgf
  \end{center}
  Dakle, svaka smislena raspodela putnika po autobusima je sirjektivna ili ``na'' funkcija (oznaka ``na''
  poti\v ce od ove slike: elementi domena se slikaju \emph{na ceo} kodomen).
\end{EX}

\bigskip

Funkcija $f : A \to B$ je \emph{bijektivna} (ka\v zemo jo\v s i da je $f$ \emph{bijekcija}) ako je injektivna i sirjektivna.
Tada funkcija $f$ preslikava razli\v cite elemente domena na razli\v cite elemente kodomena, i za svaki element kodomena
postoji ta\v cno jedan element domena koji se na njega preslikava. Bijektivna funkcija, dakle, izgleda ovako:
  \begin{center}
    \input ft08-f7.pgf
  \end{center}
svaki element skupa $B$ je ``pogo\dj en'' \emph{ta\v cno jednom} strelicom.

\begin{EX}\label{ft08.ex.DEC-SIXBIT}
  Bijektivne funkcije se tipi\v cno koriste za potrebe kodiranja. \emph{Kodiranje} je proces u kome se
  elementi nekog skupa zamenjuju elementima nekog drugog skupa, recimo 01-nizovima, u cilju sme\v stanja ili
  prenosa informacija. Jedan od starih i jednostavnih sistema kodiranja je DEC SIXBIT koga smo ve\'c pominjali
  ranije. On se koristio na samom po\v cetku ra\v cunarske ere za potrebe kodiranja simbola za 6-bitne ma\v sine.
  To je 6-bitni kod koji mo\v ze da opi\v se $2^6 = 64$ simbola prema slede\'coj tabeli:
  \begin{center}
    \begin{tabular}{r|rrrrrrrr}
             & ...000    & ...001   & ...010   & ...011   & ...100   & ...101   & ...110   & ...111  \\
      \hline
      000... & \verb*` ` & \verb`!` & \verb`"` & \verb`#` & \verb`$` & \verb`%` & \verb`&` & \verb`'`\\
      001... & \verb`(`  & \verb`)` & \verb`*` & \verb`+` & \verb`,` & \verb`-` & \verb`.` & \verb`/`\\
      010... & \verb`0`  & \verb`1` & \verb`2` & \verb`3` & \verb`4` & \verb`5` & \verb`6` & \verb`7`\\
      011... & \verb`8`  & \verb`9` & \verb`:` & \verb`;` & \verb`<` & \verb`=` & \verb`>` & \verb`?`\\
      100... & \verb`@`  & \verb`A` & \verb`B` & \verb`C` & \verb`D` & \verb`E` & \verb`F` & \verb`G`\\
      101... & \verb`H`  & \verb`I` & \verb`J` & \verb`K` & \verb`L` & \verb`M` & \verb`N` & \verb`O`\\
      110... & \verb`P`  & \verb`Q` & \verb`R` & \verb`S` & \verb`T` & \verb`U` & \verb`V` & \verb`W`\\
      111... & \verb`X`  & \verb`Y` & \verb`Z` & \verb`[` & \verb`\` & \verb`]` & \verb`^` & \verb`_`
    \end{tabular}
  \end{center}
  Na primer, kod simbola \verb`'A'` je 100001, a kod simbola \verb`'*'` je 001010.
  Obrnuto, kod 110100 korira simbol \verb`'T'`. Ovom tabelom je, dakle, uspostavljena bijekcija
  $f : A \to \{0,1\}^6$ izme\dj u alfabeta $A$ koji se sastoji od 64 simbola navedena u tabeli
  i skupa $\{0,1\}^6$.
  
  Jo\v s jedan va\v zan 6-bitni kod je \emph{Brajovo pismo} (odnosno, \emph{Brajov kod}) koga koriste osobe sa o\v ste\'cenim vidom.
  \v Sest bitova koda je raspore\dj eno kao \v sest ta\v ckica u tabelu $3 \times 2$, pri \v cemu su neke ta\v ckice
  ispup\v cene, a neke ne. Tako svaki od 01-nizova iz skupa $\{0, 1\}^6$ generi\v se druga\v ciji reljef koga osobe
  sa o\v ste\'cenim vidom mogu da osete pod prstima i tako prepoznaju o kom simbolu se radi.
\end{EX}

\begin{EX}
  Neka je $A = \calP(\{1, 2, 3, 4, 5\})$ i $B = \{0,1\}^5$.
  Svaki element $X$ skupa $A$ se mo\v ze kodirati 01-nizom $b_1 b_2 b_3 b_4 b_5 \in \{0, 1\}^5$
  gde je
  $$
    b_i = \begin{cases}
      0, & i \notin X,\\
      1, & i \in X.
    \end{cases}
  $$
  Na primer, neka je  i $X = \{1, 3, 4\}$. Tada je $X$ predstavljen nizom $10110$
  zato \v sto $1 \in B$, $2 \notin B$, $3 \in B$ itd. Jasno, je da je prazan skup predstavljen nizom $00000$,
  a ceo skup $\{1m 2m 3m 4m 5\}$ nizom $11111$. Evo tabelarnog pregleda:
  $$
    \begin{array}{c|ccccc}
                        & 1 & 2 & 3 & 4 & 5 \\ \hline
      \0                & 0 & 0 & 0 & 0 & 0 \\
      \{1, 3, 4\}       & 0 & 1 & 0 & 1 & 1 \\
      \{1, 2, 3, 4, 5\} & 1 & 1 & 1 & 1 & 1 
    \end{array}
  $$
  Lako se vidi da va\v zi i i obrnuto: svaki 01-niz odre\dj uje ta\v cno jedan element skupa $A$. Na primer
  niz $01101$ odre\dj uje skup $\{2, 3, 5\}$. Na ovaj na\v cin je, dakle, uspostavljena bijekcija
  $f : \calP(\{1, 2, 3, 4, 5\}) \to \{0,1\}^5$.
\end{EX}

\begin{EX}
  Funkcija $f : \NN \to \NN$ data sa $f(n) = n + 1$ je injektivna ($n \ne m \Rightarrow n + 1 \ne m + 1$),
  ali nije sirjektivna (jer ne postoji $n \in \NN$ takvo da je $f(n) = 1$).

  Funkcija $f : \ZZ \to \ZZ$ data sa $f(n) = n + 1$ je bijektivna (injektivnost sledi kao gore, a sirjektivnost je o\v cigledna
  jer je $f(n - 1) = n$, pa za svako $n$ postoji neko $m$ takvo da je $f(m) = n$).

  Funkcija $f : \ZZ \to \NN$ data sa $f(n) = |n + 1|$ je sirjektivna, ali nije injektivna
  (sirjektivnost sledi lako jer je $f(n - 1) = n$; funkcija nije injektivna jer $f(-5) = f(3) = 4$).
\end{EX}

\begin{EX}\label{ft08.ex.id-A}
  Za svaki skup $A$ postoji funkcija $\id_A : A \to A$ koja je definisana ovako: $\id_A(a) = a$ za sve $a \in A$.
  Ovako definisana funkcija je zove \emph{identi\v cka funkcija} i ona ne radi ni\v sta -- prosto svaki element preslika u sebe.
  \begin{center}
    \input ft08-f11.pgf
  \end{center}
  To je, o\v cito, bijekcija. Lako se pokazuje (i bi\'ce nam va\v zno kasnije) da za svaku funkciju $f : A \to B$ va\v zi:
  $$
    \id_B \circ f = f \circ \id_A = f.
  $$
\end{EX}

\begin{THM}
  \begin{enumerate}
  \item
    Kompozicija dve injektivne funkcije je injektivna funkcija.
  \item
    Kompozicija dve sirjektivne funkcije je sirjektivna funkcija.
  \item
    Kompozicija dve bijektivne funkcije je bijektivna funkcija.
  \end{enumerate}
\end{THM}
\begin{proof}[Dokaz]
  Neka su $f : A \to B$ i $g : B \to C$ proizvoljne funkcije.

  \medskip

  (1)
  Pretpostavimo da su $f$ i $g$ injektivne i poka\v zimo da je tada i njihova kompozicija $g \circ f : A \to C$ injektivna.
  Neka su $a, b \in A$ takvi da je $a \ne b$. Zato \v sto je $f$ injektivna funkcija zaklju\v cujemo da je
  $f(a) \ne f(b)$. No tada je i $g(f(a)) \ne g(f(b))$ zato \v sto je $g$ injektivna funkcija.
  \begin{center}
    \input ft08-f8.pgf
  \end{center}
  Dakle, pokazali smo
  da za proizvoljne $a, b \in A$,
  $$
    a \ne b \Rightarrow (g \circ f)(a) \ne (g \circ f)(b),
  $$
  \v sto zna\v ci da je $g \circ f$ injektivna funkcija.

  \medskip
  
  (2)
  Pretpostavimo da su $f$ i $g$ sirjektivne i poka\v zimo da je tada i njihova kompozicija $g \circ f : A \to C$ sirjektivna.
  Uzmimo proizvoljno $c \in C$. Zato \v sto je $g$ sirjektivna funkcija postoji neko $b \in B$ takvo da je $g(b) = c$.
  Zato \v sto je $f$ sirjektivna funkcija postoji neko $a \in A$ takvo da je $f(a) = b$. Dakle,
  $c = g(b) = g(f(a))$.
  \begin{center}
    \input ft08-f10.pgf
  \end{center}
  Ovim smo pokazali da za svako $c \in C$ postoji $a \in A$ takvo da je $(g \circ f)(a) = c$ \v sto
  zna\v ci da je $g \circ f$ sirjektivna funkcija.
  
  \medskip
  
  (3)
  Sledi direktno iz (1) i (2) zato \v sto je bijektivna funkcija istovremeno i injektivna i sirjektivna.
\end{proof}

\begin{THM}
  Neka su $f : A \to B$ i $g : B \to C$ proizvoljne funkcije.
  \begin{enumerate}
  \item
    Ako je $g \circ f$ injektivna, onda je $f$ injektivna.
  \item
    Ako je $g \circ f$ sirjektivna, onda je $g$ sirjektivna.
  \item
    Ako je $g \circ f$ bijektivna, onda je $f$ injektivna, a $g$ sirjektivna.
  \end{enumerate}
\end{THM}
\begin{proof}[Dokaz]
  (1)
  Da bismo pokazali da je $f$ injektivna, pokaza\'cemo da za sve $a_1, a_2 \in A$ va\v zi slede\'ce:
  $$
    f(a_1) = f(a_2) \Rightarrow a_1 = a_2.
  $$
  Neka je $f(a_1) = f(a_2)$. Onda je $g(f(a_1)) = g(f(a_2))$, odnosno $(g \circ f)(a_1) = (g \circ f)(a_2)$.
  Zato \v sto je $g \circ f$ injektivna poslednja jednakost implicira da je $a_1 = a_2$. Dakle, pokazali smo da
  $f(a_1) = f(a_2)$ implicira $a_1 = a_2$.
  
  \medskip
  
  (2)
  Uzmimo proizvoljno $c \in C$. Zato \v sto je $g \circ f$ sirjektivna postoji $a \in A$ takvo da je
  $(g \circ f)(a) = c$. No tada je $g(f(a)) = c$ pa za $b = f(a) \in B$ imamo da je $g(b) = c$.
  Dakle, pokazali smo da za svako $c \in C$ postoji $b \in B$ takvo da je $g(b) = c$.
  %\begin{center}
  %  \input ft08-f10.pgf
  %\end{center}

  \medskip
  
  (3)
  Sledi direktno iz (1) i (2) zato \v sto je bijektivna funkcija istovremeno i injektivna i sirjektivna.
\end{proof}

\clearpage

\section{Inverzna funkcija}

\noindent
\begin{minipage}{3in}
  \hspace*{1.5em}Funkcija $f : A \to B$ je \emph{invertibilna} ako postoji funkcija $g : B \to A$ takva da je
  $$
    f(a) = b \Leftrightarrow g(b) = a
  $$
  za sve $a \in A$ i sve $b \in B$. To zna\v ci: ako $f$ preslika $a$ na $b$,
  onda $g$ to $b$ ``vrati ku\'ci'' na $a$, i obrnuto.
\end{minipage}
\hfill
\begin{minipage}{2.5in}
  \begin{center}
    \input ft08-f12.pgf
  \end{center}
\end{minipage}

\medskip

\begin{EX}
  Invertibilne funkcije se tipi\v cno koriste za potrebe \v sifrovanja. Ako je $M$ skup
  poruka (\emph{messages}), a $C$ skup \v sifri (\emph{codes}) onda funkcija
  $
    \mathit{enode} : M \to C
  $
  \v sifruje poruke, a njoj odgovaraju\'ca funkcija
  $
    \mathit{decode} : C \to M
  $
  de\v sifruje k\=odove. Da bi \v sifrovanje imalo smisla mora da va\v zi
  $$
    \mathit{encode}(m) = c \Leftrightarrow \mathit{decode}(c) = m.
  $$
  Dakle, algoritam za \v sifrovanje je smislen ako va\v zi slede\'ce:
  \begin{quote}
    \v sifrovanjem poruke $m$ se dobija k\=od $c$ ako i samo ako se de\v sifrovanjem k\=oda $c$ dobija
    polazna poruka $m$.
  \end{quote}
  Jedan va\v zan kriptosistem, RSA, smo ve\'c prodiskutovali na samom po\v cetku kursa!
\end{EX}

\begin{THM}
  Funkcija $f : A \to B$ je invertibilna ako i samo ako je bijektivna.
\end{THM}
\begin{proof}[Dokaz]
  $(\Rightarrow)$
  Neka je $f : A \to B$ invertibilna funkcija i neka je $g : B \to A$ funkcija takva da je
  $$
    f(a) = b \Leftrightarrow g(b) = a
  $$
  Lako se pokazuje da je funkcija $f$ sirjektivna. Neka je $b \in B$ proizvoljno i neka je $a = g(b)$. Tada je
  $f(a) = b$, pa za svako $b \in B$ postoji $a \in A$ takvo da je $f(a) = b$.

  Dokaza\'cemo da je $f$ injektivna funkcija tako \v sto \'cemo pokazati
  $$
    f(a_1) = f(a_2) \Rightarrow a_1 = a_2.
  $$
  Neka je $f(a_1) = f(a_2) = b$ za neko $b \in B$. Zbog $f(a_1) = b$ je $g(b) = a_1$, a zbog
  $f(a_2) = b$ je $g(b) = a_2$. Dakle, $a_1 = g(b) = a_2$.

  $(\Leftarrow)$
  Neka je $f : A \to B$ bijektivna funkcija. To zna\v ci da za svako $b \in B$ postoji ta\v cno jedno
  $a \in A$ takvo da je $f(a) = b$. Zato je sa
  $$
    g(b) = a \Leftrightarrow f(a) = b
  $$
  definisana funkcija $g : B \to A$ koja po konstrukciji ima tra\v zenu osobinu.
\end{proof}

Za dalja razmatranja \'ce nam biti korisno da preformuli\v semo pojam invertibilnosti. Naredna teorema
nam pokazuje kako svojstvo invertibilnosti mo\v zemo da zapi\v semo na drugi na\v cin. Njen dokaz je veoma
jednostavan zato \v sto se, prakti\v cno, radi o prevo\dj enju sa jednog na drugi jezik.

\begin{THM}
  Funkcija $f : A \to B$ je invertibilna ako i samo ako postoji funkcija $g : B \to A$ takva da je
  $g \circ f = \id_A$ i $f \circ g = \id_B$.
\end{THM}
\begin{proof}[Dokaz]
  $(\Rightarrow)$
  Neka je $f : A \to B$ invertibilna funkcija. Onda postoji funkcija $g : B \to A$ takva da
  je $f(a) = b \Leftrightarrow g(b) = a$. Sada se lako pokazuje da je $g \circ f = \id_A$. Uzmimo proizvoljno $a \in A$
  i neka je $b = f(a)$. Tada je $g(b) = a$ pa je $(g \circ f)(a) = g(f(a)) = g(b) = a = \id_A(a)$. Zato je $g \circ f = \id_A$.
  Na isti na\v cin se pokazuje da je $f \circ g = \id_B$.
  
  $(\Leftarrow)$
  Neka je $g : B \to A$ funkcija takva da je $g \circ f = \id_A$ i $f \circ g = \id_B$. Poka\v zimo da tada va\v zi
  $$
    f(a) = b \Leftrightarrow g(b) = a.
  $$
  Neka je $f(a) = b$. Zato \v sto je $g \circ f = \id_A$ dobijamo da je
  $g(f(a)) = a$, pa kako je $f(a) = b$ sledi da je $g(b) = a$. Druga implikacija se dokazuje na isti na\v cin.
\end{proof}

\begin{THM}
  Neka je $f : A \to B$ invertibilna funkcija. Tada postoji ta\v cno jedna funkcija $g : B \to A$ takva da je
  $g \circ f = \id_A$ i $f \circ g = \id_B$.
\end{THM}
\begin{proof}[Dokaz]
  Pretpotavimo da su $g_1 : B \to A$ i $g_2 : B \to A$ funkcije sa navedenom osobinom. Doka\v zimo da je tada $g_1 = g_2$.
  Prema konstataciji iz Primera~\ref{ft08.ex.id-A} znamo da je $g_1 = g_1 \circ \id_B$, a prema pretpostavci je
  $\id_B = f \circ g_2$. Zato je
  $$
    g_1 = g_1 \circ \id_B = g_1 \circ (f \circ g_2) = (g_1 \circ f) \circ g_2 = \id_A \circ g_2 = g_2,
  $$
  gde smo jo\v s jednom koristili pretpostavku teoreme i konstatciju iz Primera~\ref{ft08.ex.id-A}.
\end{proof}

Ako je $f : A \to B$ invertibilna funkcija, onda jedinstvenu funkciju $g : B \to A$ za koju je
$g \circ f = \id_A$ i $f \circ g = \id_B$ zovemo \emph{inverzna funkcija funkcije $f$} i ozna\v cavamo je sa~$f^{-1}$.

\begin{THM}
  Neka je $f$ invertibilna funkcija. Tada je $(f^{-1})^{-1} = f$.
\end{THM}


\section[UPORE\Dj IVANJE KONA\v CNIH SKUPOVA]{Upore\dj ivanje kona\v cnih skupova}

Neka su $A$ i $B$ kona\v cni skupovi. Da bismo utvrdili da li $A$ i $B$ imaju isti broj elemenata
ili neki od njih ima vi\v se elemenata od ono drugog dovoljno je prebrojati elemente jednog i drugog skupa
i uporediti dobijene brojeve. Me\dj utim, postoje situacije u kojima nije lako utvrditi $|A|$ i $|B|$
(videti Primer~\ref{ft08.ex.22-34}), ali se ipak mo\v ze utvrditi odnos brojeva $|A|$ i $|B|$ koriste\'ci
mala lukavstva.

\paragraph{Princip bijekcije.}
  Neka su $A$ i $B$ kona\v cni skupovi. Tada je $|A| = |B|$ ako i samo ako postoji bar jedna bijektivna funkcija $f : A \to B$.
\begin{center}
  \input ft08-f7.pgf
\end{center}

\begin{EX}\label{ft08.ex.22-34}
  Kojih brojeva ima vi\v se u skupu $\{1, 2, \ldots, 999\;999\}$: onih \v ciji zbir cifara je~22, ili onih \v ciji zbir cifara je~32?
\end{EX}
\begin{proof}[Re\v senje]
  Neka je $S = \{1, 2, \ldots, 999\;999\}$. Svaki broj iz skupa $S$ mo\v zemo zapisati kao \v sectocifren broj dopisivanjem nula
  na po\v cetak ako je to potrebno. Na primer, 499 \'cemo zapisati kao 000\;499, a 910\;877 \'cemo ostaviti u tom obliku.
  Neka je
  \begin{align*}
    A &= \{ a \in S : \text{ zbir cifara broja $a$ je 22}\}, \text{ i}\\
    B &= \{ b \in S : \text{ zbir cifara broja $b$ je 32}\}.
  \end{align*}
  Pokaza\'cemo da je $|A| = |B|$ tako \v sto \'cemo konstruisati bijekciju $f : A \to B$.
  Za \v sestocifreni broj $\overline{a_1 a_2 a_3 a_4 a_5 a_6} \in A$ stavimo
  $$
    f(\overline{a_1 a_2 a_3 a_4 a_5 a_6}) = \overline{b_1 b_2 b_3 b_4 b_5 b_6}
  $$
  gde je $b_i = 9 - a_i$. Na primer, $f(000\;499) = 999\;500$.
  Da bismo pokazali da je $f$ zaista funkcija iz $A$ u $B$ treba da poka\v zemo da se svaki element skupa $A$ preslikava
  na neki element skupa $B$. Neka je \hbox{$\overline{a_1 a_2 a_3 a_4 a_5 a_6} \in A$}. Tada je
  $a_1 + a_2 + a_3 + a_4 + a_5 + a_6 = 22$. Neka je, dalje,
  $f(\overline{a_1 a_2 a_3 a_4 a_5 a_6}) = \overline{b_1 b_2 b_3 b_4 b_5 b_6}$. Tada je
  \begin{align*}
  b_1 + b_2 &+ b_3 + b_4 + b_5 + b_6 = \\
  &= (9 - a_1) + (9 - a_2) + (9 - a_3) + (9 - a_4) + (9 - a_5) + (9 - a_6)\\
  &= 54 - (a_1 + a_2 + a_3 + a_4 + a_5 + a_6) = 54 - 22 = 32.
  \end{align*}
  Dakle, $f$ je zaista funkcija koja preslikava elemente skupa $A$ u elemente skupa~$B$.
  Funkcija $f$ je injektivna po konstrukciji, a lako se pokazuje i da je sirjektivna. Neka je
  $\overline{b_1 b_2 b_3 b_4 b_5 b_6} \in B$. Napravimo \v sestocifreni broj
  $\overline{a_1 a_2 a_3 a_4 a_5 a_6}$ tako da je $a_i = 9 - b_i$. Tada se kao i gore pokazuje
  da je $\overline{a_1 a_2 a_3 a_4 a_5 a_6} \in A$, dok je
  $f(\overline{a_1 a_2 a_3 a_4 a_5 a_6}) = \overline{b_1 b_2 b_3 b_4 b_5 b_6}$
  zato \v sto je $9 - a_i = 9 - (9 - b_i) = b_i$. Na primer, ako \v zelimo da utvrdimo koji broj se preslikava
  na broj $910\;877 \in B$ samo redom oduzmemo njegove cifre od 9 i tako dobijamo da je
  $f(089\;122) = 910\;877$.
\end{proof}

\paragraph{Princip sirjekcije.}
  Neka su $A$ i $B$ kona\v cni skupovi. Tada je $|A| \ge |B|$ ako i samo ako postoji bar jedna sirjektivna funkcija $f : A \to B$.
\begin{center}
  \input ft08-f13.pgf
\end{center}

\paragraph{Princip injekcije.}
  Neka su $A$ i $B$ kona\v cni skupovi. Tada je $|A| \le |B|$ ako i samo ako postoji bar jedna injektivna funkcija $f : A \to B$.
\begin{center}
  \input ft08-f14.pgf
\end{center}

Kontrapozicija Principa injekcije je tvr\dj enje poznato kod nas kao Dirihleov princip, a na engleskom govornom podru\v cju kao
\emph{Pigeon-Hole Principle} (``princip golubova i rupa'').

\paragraph{Dirihleov princip.}
  Neka su $A$ i $B$ kona\v cni skupovi. Tada je $|A| > |B|$ ako i samo ako ne postoji injektivno preslikavanje $f : A \to B$.
  
\medskip

Poetskija formulacija ovog principa se zove \emph{Pigeon-Hole Principle} (``princip golubova i rupa'') i glasi ovako:
ako upla\v simo $n + 1$ golubova i oni odlete i sakriju se u $n$ rupa na golubarniku, onda sigurno postoji bar jedna rupa
u kojoj su se sakrila bar dva goluba.

\begin{EX}
  Vojnik je deset puta pucao u kvadratnu metu dimenzija $210\mathit{cm} \times 210\mathit{cm}$ i svaki put pogodio.
  Dokazati da postoje dva pogotka koji su na rastojanju manjem of $1\mathit{m}$.
\end{EX}

\noindent
\begin{minipage}{3in}
  \begin{proof}[Re\v senje]
    Podelimo metu na devet manjih kvadrata dimenzija $70\mathit{cm} \times 70\mathit{cm}$ kao na slici pored.
    Vojnik je pucao u metu deset puta, pa prema Dirihleovom principu
    postoji manji kvadrat u kome se nalaze dva pogotka.
    Najve\'ce rastojanje koje dve ta\v cke u kvadratu mogu da postignu je du\v zina dijagonale kvadrata, \v sto zna\v ci
    da su ona dva pogotka koja se nalaze u istom kvadratu na rastojanju koje je najvi\v se
    $$
      70\mathit{cm} \cdot \sqrt 2 = 98.9949\ldots\mathit{cm} < 100\mathit{cm}.\qedhere
    $$
  \end{proof}
\end{minipage}
\hfill
\begin{minipage}{2.5in}
  \begin{center}
    \input ft08-f15.pgf
  \end{center}
\end{minipage}


\section[UPORE\Dj IVANJE BESKONA\v CNIH SKUPOVA]{Upore\dj ivanje beskona\v cnih skupova}

Neka su $A$ i $B$ beskona\v cni skupovi. Da bismo utvrdili da li $A$ i $B$ imaju isti broj elemenata
ili neki od njih ima vi\v se elemenata od ono drugog vi\v se ne mo\v zemo da se oslanjamo na ideju da \'cemo
``prebrojati'' elemente oba skupa i onda uporediti dobijene ``brojeve''.
Postoji matemati\v cki jezik koji nam omogu\'cuje da govorimo o broju elemenata beskona\v cnog skupa ali je on
veoma komplikovan i daleko prevazilazi teorijske osnove koje smo do sada izlo\v zili.

Me\dj utim, i dalje mo\v zemo koristiti ideje koje su izlo\v zene u prethodnom odeljku:
dva skupa (kona\v cna ili beskona\v cna) mo\v zemo da uporedimo po veli\v cini tako \v sto
utvrdimo kakvi tipovi funkcija se mogu uspostaviti izme\dj u njih.
Koriste\'ci iste ideje mo\v zemo precizno da ka\v zemo kada je skup kona\v can, a kada je beskona\v can.

Skup $A$ je \emph{kona\v can} ako je svaka injektivna funkcija $f : A \to A$ ujedno i bijektivna.
Skup $A$ je \emph{beskona\v can} ako postoji injektivna funkcija $f : A \to A$ koja nije bijektivna.

\begin{EX}
  Skup $A = \{1, 2, 3, 4, 5\}$ je kona\v can zato \v sto svaka injektivna funkcija $f : A \to A$ mora da
  preslika $A$ na ceo kodomen, pa je stoga bijektivna.
  
  Skup $\NN$ je beskona\v can zato \v sto je funkcija $f : \NN \to \NN$ data sa $f(n) = n + 1$ injektivna, a nije
  bijektivna (nijedan element se ne preslikava na broj~1).
\end{EX}

\begin{THM}
  \begin{enumerate}
  \item Ako je $A$ kona\v can skup i $B \subseteq A$ onda je i $B$ kona\v can skup.
  \item Ako je $A$ beskona\v can skup i $A \subseteq C$ onda je i $C$ beskona\v can skup.
  \end{enumerate}
\end{THM}
\begin{proof}[Dokaz]
  (1)
  Neka je $A$ kona\v can skup i neka je $B \subseteq A$. Neka je $f : B \to B$ proizvoljna injektivna funkcija.
  Defini\v simo funkciju $g : A \to A$ ovako
  $$
    g(x) = \begin{cases}
      f(x), & x \in B\\
      x     & x \in A \setminus B.
    \end{cases}
  $$

\noindent
\begin{minipage}{3in}
  Tada je $g$ injektivna funkcija, pa mora biti bijektivna zato \v sto je $A$ kona\v can skup.
  Odatle lako sledi da i $f$ mora biti bijektivna funkcija. Naime, $g$ je sirjektivna \v sto zna\v ci
  da je svaki element kodomena slika nekog elementa. Kako na delu $A \setminus B$ funkcija $g$ ne radi
  ni\v sta interesantno, zaklju\v cujemo da funkcija $f$ mora da preslika $B$ na ceo $B$, pa je zato $f$ sirjektivna
  funkcija.
\end{minipage}
\hfill
\begin{minipage}{2.5in}
  \begin{center}
    \input ft08-f16.pgf
  \end{center}
\end{minipage}

\medskip

\noindent
\begin{minipage}{3in}
  \hspace*{1.5em}(2)
  Neka je $A$ beskona\v can skup i neka je $A \subseteq C$. Neka je $f : A \to A$ injektivna funkcija koja nije bijektivna.
  Defini\v simo funkciju $g : C \to C$ ovako
  $$
    g(x) = \begin{cases}
      f(x), & x \in A\\
      x     & x \in C \setminus A.
    \end{cases}
  $$
  Tada je $g$ injektivna funkcija koja nije bijektivna, i zato $C$ nije kona\v can skup.\qedhere
\end{minipage}
\hfill
\begin{minipage}{2.5in}
  \begin{center}
    \input ft08-f17.pgf
  \end{center}
\end{minipage}
\end{proof}

\begin{EX}
  Skupovi $\ZZ$, $\QQ$ i $\RR$ su beskona\v cni zato \v sto je $\NN$ beskona\v can skup i $\NN \subseteq \ZZ \subseteq \QQ \subseteq \RR$.
\end{EX}

\medskip

Skupovi $A$ i $B$ (kona\v cni ili beskona\v cni) \emph{imaju isti broj elemenata}, i to zapisujemo ovako $|A| = |B|$,
ako postoji bar jedna bijekcija $f : A \to B$.

\begin{THM}
  $|\NN| = |\ZZ|$.
\end{THM}
\begin{proof}[Dokaz]
  Pokaza\'cemo da je $|\NN| = |\ZZ|$ tako \v sto \'cemo konstruisati bijekciju $f : \NN \to \ZZ$. Tabelom bijekciju $f$
  mo\v zemo da opi\v semo ovako:
  $$
    \begin{array}{c|ccccccccccccccccc}
      n    & 1 & 2  & 3 & 4  & 5 & 6  & 7 & 8  & 9 & \ldots \\ \hline
      f(n) & 0 & -1 & 1 & -2 & 2 & -3 & 3 & -4 & 4 & \ldots
    \end{array},
  $$
  a formulom ovako:
  $$
    f(n) = \begin{cases}
      \displaystyle \frac{n-1}{2}, & n \text{\ neparno},\\[5mm]
      \displaystyle -\frac{n}{2},  & n \text{\ parno}.\qquad\qedhere
    \end{cases}
  $$
\end{proof}

Skup $A$ je \emph{prebrojivo beskona\v can} ako je $|A| = |\NN|$. Ime poti\v ce od toga \v sto je $\NN = \{1, 2, 3, \ldots\}$, i
svaki prebrojivo beskona\v can skup $A$ mo\v ze da se zapi\v se na isti na\v cin (elementi mogu da mu se \emph{nabroje}):
$$
  A = \{a_1, a_2, a_3, \ldots \}.
$$
Prethodna teoreme nam, dakle, ka\v ze da je $\ZZ$ prebrojivo beskona\v can skup jer njegovi elementi mogu da se nabroje ovako:
$$
  \ZZ = \{ 0, -1, 1, -2, 2, -3, 3, \ldots \}.
$$

Skup $A$ ima \emph{manje ili jednako} elemenata kao skup $B$, i pi\v semo $|A| \le |B|$, ako postoji injektivno
preslikavanje $f : A \to B$.

\begin{THM}
  Ako je $A \subseteq B$ onda je $|A| \le |B|$.
\end{THM}
\begin{proof}[Dokaz]
  Neka je $A \subseteq B$. Tada je preslikavanje $f : A \to B$ dato sa $f(a) = a$ za sve $a \in A$ injektivno,
  pa je $|A| \le |B|$.
\end{proof}

Ako je $|A| \le |B|$ i $|A| \ne |B|$ onda pi\v semo da je $|A| < |B|$ ili da je $|B| > |A|$, kako nam odgovara.
Za kona\v cne skupove je jasno da
\begin{itemize}
\item
  va\v zi ta\v cno jedan od ova tri odnosa: $|A| < |B|$, $|A| = |B|$ ili $|A| > |B|$; kao i da
\item
  iz $|A| \le |B|$ i $|B| \le |A|$ sledi da je $|A| = |B|$.
\end{itemize}
Za beskona\v cne skupove, me\dj utim, ovi stavovi nisu nimalo o\v cigledni. Na primer,
da bi se pokazao prvi stav treba pokazati da za svaka dva beskona\v cna skupa $A$ i $B$ ili postoji injektivno preslikavanje
$f : A \to B$, ili postoji bijektivno preslikavanje iz $g : A \to B$, ili postoji injektivno preslikvanje iz $h : B \to A$.
Da bi se pokazao drugi stav treba pokazati da iz \v cinjenice
da postoji injektivno preslikavanje $f : A \to B$ i injektivno preslikavanje $g : B \to A$ sledi da postiji
bijekcija $h : A \to B$. Sre\'com, oba ova stava su ta\v cna i za beskona\v cne skupove, ali njihovi dokazi prevazilaze
ambicije ovog teksta. Zato \'cemo samo navesti odgovaraju\'ce teoreme bez dokaza.

\begin{THM}[Trihotomija]
  Neka su $A$ i $B$ proizvoljni skupovi. Tada va\v zi ta\v cno jedan od ova tri odnosa: $|A| < |B|$, $|A| = |B|$ ili $|A| > |B|$.
\end{THM}

\begin{THM}[Kantor, Bern\v stajn, \v Sreder]
  Neka su $A$ i $B$ proizvoljni skupovi. Ako je $|A| \le |B|$ i $|B| \le |A|$, onda je $|A| = |B|$.
\end{THM}

Koriste\'ci Teoremu Kantora, Bern\v stajna i \v Sredera dokaza\'cemo da je $\QQ$ prebrojivo beskona\v can skup.

\begin{THM}
  $|\NN| = |\QQ|$.
\end{THM}
\begin{proof}[Dokaz]
  Zbog $\NN \subseteq \QQ$ je $|\NN| \le |\QQ|$. Doka\v zimo sada da je $|\QQ| \le |\NN|$. Svaki racionalan broj
  $x \in \QQ$ mo\v ze na vi\v se na\v cina da se prika\v ze u obliku
  $$
    x = \frac mn
  $$
  gde je $m \in \ZZ$ i $n \in \NN$. Na primer,
  $$
    \frac{-4}{7} = \frac{-8}{14} = \frac{-44}{77} = \ldots
  $$
  Za konstrukciju koja sledi moramo da se opredelimo za jedinstven na\v cin da
  racionalan broj prika\v zemo u obliku razlomka, pa \'cemo postupiti na slede\'ci na\v cin.
  Racionalan broj $x = 0$ \'cemo prikazati kao $\frac 01$. Ako je $x \ne 0$ onda od svih oblika
  $
    x = \frac mn
  $
  odaberimo onaj za koga je $|m|$ najmanje mogu\'ce. Na primer, u slu\v caju
  $$
    x = \frac{9}{8} = \frac{18}{16} = \frac{108}{96} = \ldots
  $$
  biramo oblik $x = \frac{9}{8}$, dok za
  $$
    x = \frac{-4}{7} = \frac{-8}{14} = \frac{-44}{77} = \ldots
  $$
  biramo oblik $x = \frac{-4}{7}$.
  
  Definisa\'cemo sada funkciju $f : \QQ \to \NN$ koja razlomke kodira prirodnim brojevima oblika $2^r \cdot 3^s \cdot 5^t$
  gde stepen broja $2$ kodira znak razlomka ($r = 1$ ako je razlomak negativan, i $r = 0$ ako je razlomak nenegativan),
  stepen broja $3$ kodira apsolutnu vrednost imenioca, a stepen broja 5 kodira brojilac razlomka. Dakle,
  $$
    f(\textstyle\frac mn) = 2^r \cdot 3^{|m|} \cdot 5^n
  $$
  gde je $r = 1$ ako je $m < 0$, a u ostalim slu\v cajevima je $r = 0$. Na primer,
  \begin{align*}
    f(\textstyle\frac 01) &= 2^0 \cdot 3^0 \cdot 5^1\\[3mm]
    f(\textstyle\frac 98) &= 2^0 \cdot 3^9 \cdot 5^8\\[3mm]
    f(\textstyle\frac{-4}{7}) &= 2^1 \cdot 3^4 \cdot 5^7
  \end{align*}
  Lako se vidi da je $f$ injektivna funkcija (kodiranje je jednozna\v cno!) tako da je $|\QQ| \le |\NN|$. Prema
  Teoremi Kantora, Bern\v stajna i \v Sredera sada zaklju\v cujemo da je $|\NN| = |\QQ|$.
\end{proof}

Svi skupovi koje smo do sada videli su bili kona\v cni ili prebrojivo beskona\v cni. No, to nije kraj pri\v ce!
Kantor je pokazao da skup $\RR$ nije prebrojivo beskona\v can. On, dakle, ima vi\v se elemenata od skupa $\NN$:

\begin{THM}[Kantor]
  $|\NN| < |\RR|$.
\end{THM}
\begin{proof}[Dokaz]
  Podsetimo se, prvo, da svaki realan broj $x \in \RR$ mo\v ze da se predstavi u obliku
  $$
    x = k_0,k_1 k_2 k_3 \ldots
  $$
  gde je $k_0 \in \ZZ$ ceo deo broja $x$, a $k_1$, $k_2$, $k_3$, \ldots su decimale broja $x$. Na primer za
  $$
    x = 3,1415926\ldots
  $$
  je $k_0 = 3$, $k_1 = 1$, $k_2 = 4$, $k_3 = 1$ itd, dok za
  $$
    x = -12,5
  $$
  imamo da je $k_0 = -12$, $k_1 = 5$, $k_2 = 0$, $k_3 = 0$ itd.

  Pre\dj imo sada na dokaz teoreme.
  Zbog $\NN \subseteq \RR$ je $|\NN| \le |\RR|$. Da bismo dokazali teoremu dovoljno je da poka\v zemo da $|\NN| \ne |\RR|$.
  Pretpostavimo suprotno: $|\NN| = |\RR|$. Tada postoji bijekcija $f : \NN \to \RR$. Neka je
  \begin{align*}
    f(1) &= a_0, a_1 a_2 a_3 a_4 \ldots\\
    f(2) &= b_0, b_1 b_2 b_3 b_4 \ldots\\
    f(3) &= c_0, c_1 c_2 c_3 c_4 \ldots\\
         &\vdots
  \end{align*}
  Posmatrajmo broj $\gamma \in \RR$ koji je konstruisan ovako:
  $$
    \gamma = 0, z_1 z_2 z_3 z_4 \ldots
  $$
  gde je
  $$
    z_1 = \begin{cases}
      1, & a_1 = 0\\
      0, & a_1 \ne 0,
    \end{cases},
    \quad
    z_2 = \begin{cases}
      1, & b_2 = 0\\
      0, & b_2 \ne 0,
    \end{cases},
    \quad
    z_3 = \begin{cases}
      1, & c_3 = 0\\
      0, & c_3 \ne 0,
    \end{cases},
    \quad
    \text{itd.}
  $$
  Broj $\gamma$ je konstruisan tako da bude $z_1 \ne a_1$, $z_2 \ne b_2$, $z_3 \ne c_3$ itd.
  Dakle, uo\v cili smo dijagonalu me\dj u decimalama niza
  \begin{align*}
    f(1) &= a_0, \underline{a_1} a_2 a_3 a_4 \ldots\\
    f(2) &= b_0, b_1 \underline{b_2} b_3 b_4 \ldots\\
    f(3) &= c_0, c_1 c_2 \underline{c_3} c_4 \ldots\\
         &\vdots
  \end{align*}
  i konstruisali broj $\gamma$ tako da je $n$-ta decimala broja $\gamma$ razli\v cita od $n$-te decimale broja $f(n)$.
  Zbog toga se ova konstrukcija
  zove \emph{Kantorov dijagonalni postupak} ili \emph{Kantorova dijagonalizacija}.
  Po\v sto je $f : \NN \to \RR$ bijekcija i po\v sto je $\gamma \in \RR$ postoji neko $n$ takvo da je
  $f(n) = \gamma$. Neka je
  $$
    f(n) = k_0, k_1 k_2 \ldots \underline{k_n} \ldots
  $$
  Tada je
  $$
    k_0, k_1 k_2 \ldots \underline{k_n} \ldots = 0, z_1 z_2 \ldots z_n \ldots
  $$
  odakle sledi da je $k_n = z_n$. To je u kontradikciji sa konstrukcijom broja $\gamma$ koja garantuje da je
  $n$-ta decimala broja $\gamma$ razli\v cita od $n$-te decimale broja $f(n)$.
\end{proof}

Videli smo svojevremeno (Teorema~\ref{ft06.thm.A-PA}) da za kona\v can skup $A$ sa $n$ elemenata
skup $\calP(A)$ ima vi\v se elemenata od skupa $A$. Sada \'cemo pokazati da je to ta\v cno i za
beskona\v cne skupove.

\begin{THM}
  Za svaki (kona\v can ili beskona\v can) skup $A$ je $|A| < |\calP(A)|$.
\end{THM}
\begin{proof}[Dokaz]
  Poka\v zimo, prvo, da je $|A| \le |\calP(A)|$. Neka je $f : A \to \calP(A)$ funkcija data sa
  $f(x) = \{x\}$. Na primer, za $A = \{1, 2, 3\}$ je $\calP(A) = \{\0$, $\{1\}$, $\{2\}$, $\{3\}$, $\{1,2\}$,
  $\{1,3\}$, $\{2,3\}$, $\{1,2,3\}\}$,
  a funkcija $f$ je definisana tako da je $f(1) = \{1\}$, $f(2) = \{2\}$ i $f(3) = \{3\}$.
  Jasno je da je $f$ injektivna funkcija, i zato je $|A| \le |\calP(A)|$.

  Da bismo dokazali teoremu dovoljno je da poka\v zemo da $|A| \ne |\calP(A)|$.
  Pretpostavimo suprotno: $|A| = |\calP(A)|$. Tada postoji bijekcija $g : A \to \calP(A)$. Neka je
  $$
    B = \{ x \in A : x \notin g(x) \}.
  $$
  O\v cito je $B \in \calP(A)$, pa kako je $g$ bijekcija, postoji neko $z$ takvo da je $g(z) = B$.
  Pogledajmo sada da li je $z \in B$. Ako je $z \in B$ onda $z \notin g(z) = B$. S druge strane,
  ako $z \notin B$ onda $z \notin g(z)$, pa je $z \in B$ jer je skup $B$ tako definisan.
  Dakle, $z \in B \Leftrightarrow z \notin B$.
  Kontradikcija.
\end{proof}

Ve\'c smo videli da je $|\NN| < |\RR|$, a prethodna teorema nam ka\v ze da je
\hbox{$|\RR| < |\calP(\RR)|$}. Dakle, imamo niz beskona\v cnih skupova od kojih svaki ima
ve\'ci broj elemenata od onog prethodnog:
$$
  |\NN| < |\RR| < |\calP(\RR)| < |\calP(\calP(\RR))| < |\calP(\calP(\calP(\RR)))| < \ldots
$$
Situacija je, zapravo, mnogo interesantnija, ali je to po\v cetak jedne sasvim druge pri\v ce.







